\documentclass[11pt]{article}

\usepackage[hmargin=0.86in,vmargin=0.72in]{geometry}
\usepackage[T1]{fontenc}
\usepackage{lmodern}
\usepackage{microtype}
\usepackage{amsmath,amssymb,amsthm,mathtools}
\usepackage{booktabs,tabularx,array,longtable}
\usepackage{graphicx}
\usepackage{caption}
\usepackage{float}
\usepackage{enumitem}
\usepackage{xcolor}
\usepackage[authoryear,round]{natbib}
\usepackage[hidelinks]{hyperref}

\setlength{\parindent}{0pt}
\setlength{\parskip}{4pt}
\setlength{\emergencystretch}{3em}
\allowdisplaybreaks
\setlist[enumerate]{leftmargin=1.45em,itemsep=2pt,topsep=3pt}
\setlist[itemize]{leftmargin=1.45em,itemsep=2pt,topsep=3pt}
\captionsetup{font=small,labelfont=bf,skip=4pt}

\newtheorem{proposition}{Proposition}
\newtheorem{lemma}{Lemma}
\theoremstyle{definition}
\newtheorem{definition}{Definition}
\newtheorem{assumption}{Assumption}

\newcommand{\E}{\mathbb E}
\newcommand{\1}{\mathbf 1}
\newcommand{\X}{\mathcal X}
\newcommand{\T}{\mathcal T}
\newcommand{\Bcal}{\mathcal B}

\title{Intergenerational Housing, Fertility, and Demographic Transition\\
\large Theory and quantitative architecture}
\author{Tommaso De Santo}
\date{August 2026}

\begin{document}
\maketitle

\begin{abstract}
Young households need additional space when they have children, but access to
family-sized housing depends on tenure, liquid wealth, and the housing retained
by older owners. I study the resulting feedback between fertility and the
housing market. A two-generation model shows that reproductive stationarity
selects the long-run housing cost, while market clearing selects population
scale. It also gives exact conditions under which reproduction can be below
replacement while household population and current housing costs continue to rise: a
large inherited cohort must add more housing demand as it ages than is lost
from the smaller cohort behind it. I then embed the mechanism in a lifecycle
model with sequential births, income risk, saving, tenure choice, collateral
constraints, child maturation, survival, and bequests. The calendar-time
equilibrium propagates the unnormalized joint distribution and clears the
housing market along the transition. A reproducible numerical illustration
separates inherited demographic momentum from a no-fertility-decline path and
compares policies that relax young access with policies that release housing
retained by older households. Finally, an audit of the maintained quantitative
specification shows what remains before the transition can be interpreted as a
U.S. quantitative result: the existing fertility block is below replacement
and responds only weakly to housing costs. The framework is therefore a
coherent theoretical and computational design; estimating the U.S. transition
remains a separate empirical and numerical task.
\end{abstract}

\section{The mechanism}

Children raise the demand for space at the same ages when households have low
liquid wealth and limited access to ownership. Large rental units are scarce,
and buying a family-sized home requires a down payment. Older households, by
contrast, own much of the large-housing stock and may retain it after their
children leave because moving is costly and housing is part of the estate they
plan to leave. The allocation of housing across generations can therefore
raise the effective cost of a child even when the aggregate stock is large.

The feedback has two directions. First, a higher housing cost or a tighter
access constraint reduces fertility. Second, children eventually form the next
generation of households, so fertility changes future housing demand. The
first direction is a household mechanism; the second is a demographic
equilibrium condition. Combining them changes both the transition and the
long-run incidence of housing policy.

The analysis proceeds in two layers. Section~\ref{sec:two_generation} makes
the equilibrium logic transparent with young and old households. Sections
\ref{sec:lifecycle}--\ref{sec:calendar_equilibrium} retain the paper's
lifecycle, tenure, and intergenerational structure. The numerical section then
uses an age-structured version of the same accounting to show the transition.
The maintained quantitative model is evaluated separately, so illustrative
results are not confused with calibrated estimates.

\section{Two-generation theory}
\label{sec:two_generation}

\subsection{Environment}

\textbf{Young households.} At date $t$, a mass $Y_t$ of young households has
income $y$. A young household chooses consumption $c$, adult housing $h$, and
children $n$. Preferences are
\begin{equation}
 u(c,h,n)=c+\alpha\log h+\theta\log n,
 \qquad \alpha>0,\quad \theta>0.
 \label{eq:toy_utility}
\end{equation}
The parameter $\theta$ is the value of children. Quasi-linearity removes income
effects and keeps the price mechanism available in closed form.
Throughout this section, income is high enough that the interior allocation is
feasible: $y>\alpha+\theta$.

\textbf{Child costs and young access.} Each child requires $q\geq0$ units of
ordinary goods and $a>0$ units of housing. The market cost of that housing is
$ap$, where $p$ is the price of housing services. A goods-equivalent wedge
$\omega\geq0$ captures the additional cost of obtaining family space through a
small rental market, a down payment, or another tenure constraint. The full
cost of a child is therefore $q+a(p+\omega)$.

\textbf{Old households.} A mass $O_t$ of old households no longer chooses
fertility. After solving its retention and estate problem, an old household's
housing demand is summarized by
\begin{equation}
 d_O(p;\tau)=\frac{\varphi(\tau)}{p},
 \qquad \varphi'(\tau)>0.
 \label{eq:old_demand}
\end{equation}
The parameter $\tau$ indexes forces that keep old households in housing,
including transaction costs and the estate motive. The lifecycle model below
derives retention from choices rather than imposing equation
\eqref{eq:old_demand}.

\textbf{Housing.} Aggregate housing services are supplied according to
\begin{equation}
 H^s(p)=\bar H p^\varepsilon,
 \qquad \bar H>0,\quad \varepsilon\geq0.
 \label{eq:toy_supply}
\end{equation}
The fixed-stock case is $\varepsilon=0$. The number $\bar n>0$ is effective
replacement fertility: after survival and household formation, a young
household with $\bar n$ children produces one young household in the next
generation.

In this static section, $p$ is the current cost of one unit of housing
services. The lifecycle model distinguishes that current rent, denoted $r_t$,
from the asset price $P_t$ that capitalizes future rents. They move
proportionally only under a constant user cost and a stationary price path.

\subsection{The young household}

Taking $p$ as given, a young household solves
\begin{equation}
 \max_{c,h,n}\;c+\alpha\log h+\theta\log n
 \quad\text{subject to}\quad
 c+ph+\bigl[q+a(p+\omega)\bigr]n=y.
 \label{eq:young_problem}
\end{equation}
At an interior solution, child demand and total housing demand are
\begin{equation}
 n_Y(p;\theta,\omega)=\frac{\theta}{q+a(p+\omega)},
 \qquad
 d_Y(p;\theta,\omega)=\frac{\alpha}{p}+a n_Y(p;\theta,\omega).
 \label{eq:young_demands}
\end{equation}
A higher market price and a tighter access wedge have the same sign in the
child decision. They differ in equilibrium incidence: the market price clears
housing, while the access wedge changes which households can obtain it.

\subsection{Equilibrium and the steady state}

Housing clearing adds young and old demand:
\begin{equation}
 \bar H p_t^\varepsilon
 =Y_t d_Y(p_t;\theta,\omega)+O_t d_O(p_t;\tau).
 \label{eq:two_generation_clearing}
\end{equation}
For inherited cohort masses, the left side rises and the right side falls with
$p_t$, so the clearing price is unique.

Children become the next young cohort, while the current young become old:
\begin{equation}
 Y_{t+1}=\mathcal R(p_t;\theta,\omega)Y_t,
 \qquad O_{t+1}=Y_t,
 \qquad
 \mathcal R(p;\theta,\omega)=\frac{n_Y(p;\theta,\omega)}{\bar n}.
 \label{eq:two_generation_law}
\end{equation}
The function $\mathcal R$ is effective reproduction. It is strictly decreasing
in the housing cost and the young-access wedge.

\begin{definition}[Reproductive steady state]
A reproductive steady state is a price $p^*>0$, positive cohort masses
$(Y^*,O^*)$, and household choices such that households optimize, housing
clears, and $(Y_{t+1},O_{t+1})=(Y_t,O_t)$.
\end{definition}

\begin{proposition}[Steady state and policy incidence]
\label{prop:toy_ss}
A unique positive reproductive steady state exists if and only if
$\theta>\bar n(q+a\omega)$. It is
\begin{equation}
 p^*=\frac{\theta/\bar n-q}{a}-\omega,
 \qquad
 Y^*=O^*=\frac{\bar H(p^*)^{\varepsilon+1}}
 {\alpha+\varphi(\tau)+a\bar n p^*}.
 \label{eq:toy_ss}
\end{equation}
A lower value of children $\theta$ or a tighter young-access wedge $\omega$ lowers
both the long-run housing cost and household population. Stronger old retention
raises $\varphi(\tau)$ and lowers population without changing the reproductive
housing cost. A proportional increase in housing supply $\bar H$ raises population in the
same proportion and leaves the reproductive housing cost unchanged.
\end{proposition}

Reproduction and housing clearing perform different jobs. Replacement selects
the housing cost at which a cohort reproduces itself. Housing clearing then
selects how many young and old households can occupy the market at that cost.
This triangular ordering is why permanent housing supply is capitalized into
population rather than the long-run price in the closed economy.

The existence condition is economically sharp. Even if market housing were
free, a child would still cost $q+a\omega$. If $\theta\leq\bar n(q+a\omega)$, cheaper
market housing alone cannot restore replacement.

\begin{figure}[ht!]
\centering
\includegraphics[width=0.98\textwidth]{../output/model/fertility_population_housing_transition_note/steady_state_comparative_statics.pdf}
\caption{Reproduction selects the steady-state housing cost, and housing
clearing maps that cost into population. After a fall in the value of
children, fertility is initially below replacement. A lower population then
reduces the price and partly offsets the original shock.}
\label{fig:toy_feedback}
\end{figure}

\subsection{Demographic momentum}

The steady-state result does not imply that population and prices fall
immediately. Define the old-to-young ratio $\mu_t=O_t/Y_t$ and current
effective reproduction $\mathcal R_t=\mathcal R(p_t;\theta,\omega)$. The next result compares the
smaller young cohort created by low fertility with the inherited cohort that
is currently moving into old age.

\begin{proposition}[Low fertility with rising population and housing costs]
\label{prop:toy_momentum}
Suppose the preference, access, retention, and supply primitives are unchanged
between $t$ and $t+1$, and $\mathcal R_t<1$. Household population rises between $t$ and
$t+1$ if and only if
\begin{equation}
 \mu_t<\mathcal R_t<1.
 \label{eq:population_momentum}
\end{equation}
The housing cost rises if and only if
\begin{equation}
 (1-\mu_t)d_O(p_t;\tau)
 >(1-\mathcal R_t)d_Y(p_t;\theta,\omega).
 \label{eq:price_momentum}
\end{equation}
Hence low fertility, population growth, and housing-cost growth coexist exactly when
\begin{equation}
 \mu_t<\min\left\{\mathcal R_t,
 1-(1-\mathcal R_t)\frac{d_Y(p_t;\theta,\omega)}{d_O(p_t;\tau)}\right\}.
 \label{eq:joint_momentum}
\end{equation}
\end{proposition}

Equation~\eqref{eq:price_momentum} is an accounting test. The smaller next
young cohort removes $(1-\mathcal R_t)Y_t d_Y$ units of demand. The unusually large
current young cohort adds $(Y_t-O_t)d_O$ units as it ages. The housing cost rises when the
second force is larger. From an exact stationary age distribution,
$\mu_t=1$, so both population and the housing cost fall after a fertility decline. A
positive transition requires an inherited cohort bulge generated by earlier
fertility, immigration, or both.

The fertility decline nevertheless lowers housing costs relative to the no-shock
path. Demographic momentum can make the realized cost rise over time; it
cannot reverse the counterfactual effect of the shock.

\subsection{Intergenerational policy}

Proposition~\ref{prop:toy_ss} distinguishes three interventions. Relaxing the
young-access wedge $\omega$ raises fertility at a given cost and supports a higher
reproductive housing cost and population. Reducing old retention $\tau$ releases
housing, lowers the transition cost, and raises fertility during adjustment;
in the new closed-economy steady state it raises population while the
reproductive housing cost returns to the level selected by $\theta$ and $\omega$. Expanding
housing supply has the same long-run housing-cost neutrality but does not directly
change which generation occupies the stock.

This is where the demographic closure strengthens the paper's original
intergenerational mechanism. Reallocating one home toward a constrained young
household can affect its current tenure and fertility, and that birth changes
the size of the future entrant cohort. The second effect is absent when the
adult population is normalized after every solution.

\section{Lifecycle environment}
\label{sec:lifecycle}

The quantitative model retains the mechanism but replaces the representative
young and old households with finite lives, sequential births, income risk,
saving, discrete ownership, continuous rental housing, survival, and bequests.
It is a one-market model: the focus is intergenerational allocation and tenure,
not spatial sorting across locations.

\subsection{States and timing}

One model period is four years. Thus the household problem uses
$\beta=\beta_{\rm annual}^{4}$ and a four-year gross liquid return $R_b$;
maintenance, property taxes, and other flow rates are converted to the same
period before entering a budget. A household enters adult life at age 18 and
can live through age 85. Its beginning-of-period state is
\begin{equation}
 x=(a,b,\xi,o,n,\mathbf c)\in\X.
 \label{eq:full_state}
\end{equation}
Here $a$ is age, $b$ is liquid net wealth, $\xi$ is persistent earnings, $o$ is
the inherited tenure and owner-size state, $n\in\{0,1,2,3+\}$ is children ever
born, and $\mathbf c$ is the vector of dependent-child ages or maturation
stages. Let $m(\mathbf c)$ be the number currently at home. Tracking the stages
separately prevents a later birth from resetting the maturation clock of an
older sibling. The retained E5b code instead uses one shared child clock, so
$\mathbf c$ is a proposed timing repair, not merely a new population closure.
It changes entrant timing and requires re-estimation before quantitative use.
Conditional on household continuation, each child's stage
advances independently. The separate population pipeline introduced below
preserves the child's future-entry claim even if the parental household exits.

The within-period order is: the household observes fertility shocks and, if
eligible, chooses whether to attempt the next birth; it chooses whether to
rent or own and, if owning, a house size; it chooses consumption, rental space,
and next-period liquid wealth; income and child states then evolve, survival is
realized, and surviving households age. A successful birth increases parity by
one. Children who mature generate potential adult entrants rather than being
added mechanically to their parent's household state.

\subsection{Preferences, children, and bequests}

Let $m=m(\mathbf c)$, let $e(m)$ be an equivalence scale, and let $\alpha(m)$ be
the consumption share when $m$ children are at home. The maintained sequential
specification uses
\begin{equation}
 e(m)=\left(\frac{2+0.7m}{2}\right)^{0.7},
 \qquad
 \alpha(m)=\alpha_0-\Delta_0\1\{m>0\}-\Delta_1 m.
 \label{eq:eqscale}
\end{equation}
The first object raises total household needs; the second shifts expenditure
toward housing when children are present. For consumption $c$ and housing
services $s$, flow utility is
\begin{equation}
 u(c,s;m)
 =e(m)^{\sigma-1}
 \frac{\left[c^{\alpha(m)}s^{1-\alpha(m)}\right]^{1-\sigma}}
 {1-\sigma}
 +\psi m,
 \qquad \sigma>1.
 \label{eq:full_utility}
\end{equation}
Rental services equal the chosen quantity $h^R$. An owner of physical size
$H_j$ receives services $\chi H_j$, where $\chi$ is the owner-service premium.
Physical rooms, not utility services, enter market clearing.

At death, liquid wealth and net sale proceeds from housing form estate wealth
$w^e$. The maintained child-invariant warm glow is
\begin{equation}
 \Bcal(w^e)
 =\theta_0\frac{(\theta_1+\max\{w^e,0\})^{1-\sigma}}
 {1-\sigma}.
 \label{eq:full_bequest}
\end{equation}
Bequests affect saving and old-age housing retention. They do not determine a
particular descendant's entry wealth; the model is intergenerational in
housing allocation and cohort accounting, not fully dynastic. The displayed
utility is deliberately not recentered at zero wealth because the maintained
benchmark is not normalized; recentering would change fertility value levels
and would require re-estimation.

\subsection{Income, tenure, and housing finance}

Working income is $y(a,\xi)$, with a lifecycle profile and a persistent Markov
component. Retirees receive a pension. Liquid wealth earns gross return $R_b$.
Renters choose $h^R\in[0,\bar h^R]$. Owners choose from the physical menu
$\{H_1,\ldots,H_J\}$.

Let $P_t$ be the asset price of one unit of housing and $r_t$ its rent. A seller
pays proportional cost $\psi^s$. If the household changes from inherited
owner size $H(o)$ to new size $H(o')$, the transaction changes liquid
resources by
\begin{equation}
 \mathcal T_t(o,o')=
 \1\{H(o')\neq H(o)\}
 \left[(1-\psi^s)P_tH(o)-P_tH(o')\right].
 \label{eq:housing_transaction}
\end{equation}
Here $H(R)=0$. Keeping the same house creates no transaction. A new buyer must
finance the down payment from cash:
\begin{equation}
 \text{cash before purchase}\geq(1-\phi)P_tH(o'),
 \label{eq:downpayment}
\end{equation}
where $\phi$ is the financed share.

The maintained post-purchase rule distinguishes secured capacity from an
age-tapered unsecured position. Let
\begin{equation}
 \underline b_t^C(o')=-\phi P_tH(o'),\qquad
 \mathcal U_{a+1}(u)=min\{s^d_{a+1}\min(u,0),-D_{a+1}\},
 \label{eq:debt_rollover}
\end{equation}
where $s^d_{a+1}$ tapers existing unsecured debt toward zero and $D_{a+1}$ is
the age-specific credit line. A renter satisfies
$b'\geq\mathcal U_{a+1}(b)$; an owner satisfies
\begin{equation}
 b'\geq \underline b_t^C(o')+
 \mathcal U_{a+1}\!\left(b-\underline b_t^C(o')\right).
 \label{eq:owner_debt_floor}
\end{equation}
At the retained benchmark $D_{a+1}=0$: existing unsecured debt can initially
roll over but cannot increase, and the taper forces repayment later in life.
This is the rule in the active solver; it is not a special no-deleveraging
exception for incumbent owners.

Conditional on renting, the continuous choices satisfy
\begin{equation}
 c+r_t h^R+b'=R_b b+y(a,\xi)+\mathcal T_t(o,R)+T_t,
 \qquad 0\leq h^R\leq\bar h^R.
 \label{eq:renter_budget}
\end{equation}
Conditional on owner size $H(o')>0$, they satisfy
\begin{equation}
 c+(\delta_H+\tau_t^p)P_tH(o')+b'
 =R_b b+y(a,\xi)+\mathcal T_t(o,o')+T_t.
 \label{eq:owner_budget}
\end{equation}
Mortgage interest is embedded in the return on negative liquid wealth. The
transfer $T_t$ is zero in the baseline and is determined by the government
budget in funded policy experiments.

\subsection{Sequential fertility}

During the fertile ages, a household at parity $n<3+$ chooses between waiting
and attempting the next birth. A conception probability $f_a$ maps an attempt
into a birth. Additive Type-I extreme-value shocks have scale $\kappa_1$ on the
first-birth margin and $\kappa_C$ on continuation margins. Separating the two
scales allows the model to match childlessness and completed fertility without
interpreting a rooms response as a second-birth hazard.
Each scale-specific shock is recentered to have mean zero. This normalization
matters because $\kappa_1$ and $\kappa_C$ differ.

Conditional on a price path, let $W_t(x)$ be the optimized value of the rest of
the period conditional on entering it in state $x$. The value of waiting is
$V_t^0(x)=W_t(x)$ and the value of attempting is
$V_t^1(x)=f_aW_t(x^+)+(1-f_a)W_t(x)$, where $x^+$ increments parity and adds a
newborn child stage. The attempt probability is
\begin{equation}
 \pi_t^F(x)=
 \frac{\exp\{V_t^1(x)/\kappa(n)\}}
 {\exp\{V_t^0(x)/\kappa(n)\}+\exp\{V_t^1(x)/\kappa(n)\}},
 \qquad
 \kappa(n)=
 \begin{cases}\kappa_1,&n=0,\\ \kappa_C,&n\geq1.\end{cases}
 \label{eq:fertility_logit}
\end{equation}
Under that mean-zero normalization, the beginning-of-period inclusive value is
\begin{equation}
 V_t(x)=\kappa(n)\log\!\left[
 e^{V_t^0(x)/\kappa(n)}+e^{V_t^1(x)/\kappa(n)}\right].
 \label{eq:fertility_inclusive_value}
\end{equation}
Fertility is therefore a flow hazard; children ever born is a stock. The
calendar-time shifter used below changes the flow value of children and is not
described as an identified causal preference shock.

\subsection{The household problem}

Let $Q_t^\xi$ denote the earnings transition and $s_a$ the probability of
surviving to the next age. Let $\epsilon_t^H(o')$ be the tenure and owner-rung
taste shock. Given the full future sequence of rents, asset prices, taxes,
transfers, and child-demand shifters, the post-fertility value used above is
\begin{equation}
 W_t(x)=\E_{\epsilon_t^H}\left[
 \max_{o',c,h^R,b'}
 \left\{
 u_t(c,s;m)+\epsilon_t^H(o')+\beta s_a\E_t[V_{t+1}(x')]
 +\beta(1-s_a)\Bcal(w^e_{t+1})
 \right\}
 \right],
 \label{eq:bellman_path}
\end{equation}
subject to the appropriate budget, transaction, debt, and housing-menu
constraints above. Its continuation expectation integrates earnings and
independent child-stage transitions; conception is already integrated in
$V_t^1$. Equations~\eqref{eq:fertility_logit}--
\eqref{eq:fertility_inclusive_value} then integrate the fertility taste shock.

Equation~\eqref{eq:bellman_path} is deliberately written as a path problem.
Along a transition, current purchases occur at $P_t$ and future sales occur at
future prices. Solving a stationary Bellman problem separately at each date
would omit expected capital gains and is not a transition equilibrium.

\section{Calendar-time equilibrium}
\label{sec:calendar_equilibrium}

\subsection{Housing assets and construction}

Housing is durable. Let $K_t$ be the occupied housing stock and let $I_t$ be
net additions: gross construction minus physical retirements. Its law is
\begin{equation}
 K_{t+1}=K_t+I_t,
 \qquad I_t=\mathcal I(P_t,K_t;\zeta_t^S),
 \label{eq:housing_stock_law}
\end{equation}
where $\zeta_t^S$ is a supply shifter. A fixed-stock transition sets $I_t=0$.
The numerical diagnostic below instead uses a static,
price-elastic contemporaneous supply schedule and does not solve this stock law.
$I_t$ is a reduced-form net change in available stock. By convention there is
no unit-level depreciation risk for a maintained incumbent house: it remains
intact and can be resold.

Let $\delta_H P_t$ be maintenance that preserves an occupied unit. Competitive
ownership of a rental unit then implies the deterministic asset-pricing condition
\begin{equation}
 (R_b+\delta_H+\tau_t^p)P_t=r_t+P_{t+1}.
 \label{eq:dynamic_user_cost}
\end{equation}
When the real asset price is constant, this reduces to
$r=(R_b-1+\delta_H+\tau^p)P$, the maintained stationary user-cost relation.
The full resale value follows from the no-unit-depreciation convention above.
Transaction costs belong in household choices, not in the return condition of
a frictionless marginal landlord.

\subsection{The unnormalized cohort law}

Let $G_t$ be the finite measure of adult households over $\X$. Conditional on
the equilibrium price path, household policies define a sub-Markov household
continuation kernel $Q_t(x,A)$, expected births $f_t(x)$, and current physical
housing demand $h_t(x)$. Adult household continuation and child maturation are
different processes, so the population closure also carries a child pipeline.

Let $C_{\ell t}$ be the mass of future adult-household equivalents that is
$\ell$ periods from maturity, for $\ell=1,\ldots,L$. The conversion $\nu$
maps one birth into future household equivalents and $\varsigma_\ell$ is child
survival through stage $\ell$. With births occurring at the end of a period,
\begin{equation}
 \begin{split}
 F_t&=\int_{\X}f_t(x)\,dG_t(x),
 \qquad B_{t+1}=\varsigma_1 C_{1t},\\
 C_{\ell,t+1}&=\varsigma_{\ell+1}C_{\ell+1,t},
 \quad \ell=1,\ldots,L-1,
 \qquad C_{L,t+1}=\nu F_t.
 \end{split}
 \label{eq:child_pipeline}
\end{equation}
Thus $B_{t+1}$ is the locally born household-formation pool. Because the
pipeline is separate from $G_t$, a parental household exit does not erase a
child's future-entry claim. The vector $\mathbf c$ still determines the
parent's current needs; equation~\eqref{eq:child_pipeline} handles aggregate
demographic renewal. This aggregate closure assumes parental-household exit
does not alter child survival or eventual household formation; modeling the
child's interim guardian and housing demand would require an additional state.

Let $\mathcal M^E_{t+1}$ be an exogenous finite measure of outside entrants,
let $\rho_t$ be the fraction of locally born potential entrants retained in
the market, and let $\Gamma^B_{t+1}$ be their normalized entry-state
distribution. Actual young-adult entry is
\begin{equation}
 e_{t+1}=M_{t+1}+\rho_tB_{t+1},
 \qquad M_{t+1}=\mathcal M^E_{t+1}(\X).
 \label{eq:actual_entry}
\end{equation}
The exact calendar law is
\begin{equation}
 G_{t+1}(A)=
 \int_{\X}Q_t(x,A)\,dG_t(x)
 +\mathcal M^E_{t+1}(A)
 +\rho_tB_{t+1}\Gamma^B_{t+1}(A).
 \label{eq:level_distribution_law}
\end{equation}
This equation carries cohort size, wealth, tenure, and children through time.
It does not renormalize $G_{t+1}$. The current stationary forward routine
constructs a whole normalized cross-section from an injected entrant cohort;
a transition solver instead needs the one-period pushforward in
equation~\eqref{eq:level_distribution_law}.

Adult household population is
\begin{equation}
 S_t=\int_{\X}dG_t(x).
 \label{eq:population_and_births}
\end{equation}
Throughout the equilibrium analysis, ``population'' means adult-household
mass. Total persons require a separate aggregation of adults and dependent
children; the two series need not peak at the same date.
In the baseline transition, $\{\mathcal M^E_t,\rho_t,\Gamma^B_t\}$ are
specified exogenous paths. Endogenizing migration would require an outside
value and a location-choice equation; it is an extension, not a missing
condition in the baseline. For a national closed experiment, set
$\mathcal M^E=0$ and $\rho=1$. For an
empirical U.S. transition, immigrants should be distributed across entry ages
rather than compressed into age 18 unless the compression is stated
explicitly.

\subsection{Market clearing}

Let $h_t(x)$ be expected physical housing occupied after the current tenure
and rental choices. Housing services clear at every date:
\begin{equation}
 \int_{\X}h_t(x)\,dG_t(x)=K_t.
 \label{eq:dynamic_housing_clearing}
\end{equation}
Together, equations \eqref{eq:bellman_path},
\eqref{eq:housing_stock_law}, \eqref{eq:dynamic_user_cost},
\eqref{eq:level_distribution_law}, and \eqref{eq:dynamic_housing_clearing}
close the calendar-time model.

\begin{definition}[Transition equilibrium]
Given an initial household distribution $G_0$, child pipeline $C_0$, housing
stock $K_0$, paths for outside entrant measures, local retention, locally born
entry-state distributions, other exogenous shifters, and a terminal condition, a transition
equilibrium is a sequence of values and household policies, distributions
$\{G_t,C_t\}$, housing stocks and construction $\{K_t,I_t\}$, rents and asset
prices $\{r_t,P_t\}$, and fiscal transfers such that households optimize,
equations \eqref{eq:housing_stock_law}, \eqref{eq:dynamic_user_cost}, and
\eqref{eq:child_pipeline}--\eqref{eq:level_distribution_law} hold, housing clears by
\eqref{eq:dynamic_housing_clearing}, and every policy budget balances.
\end{definition}

The transition begins from the observed joint distribution, not an exact
steady state. Matching only the initial total population is insufficient:
demographic momentum depends on the joint age and tenure composition, while
the young-access mechanism also depends on wealth and family status.

\subsection{Stationary endpoints}

The stationary endpoint can be written without ambiguity about normalization.
At a constant price $P$, let $Q_P^\star$ be the household-continuation operator
on measures and let $\gamma_P$ be the normalized state distribution of one
entrant cohort.
Because adult life is finite, the lifetime adult measure produced by one new
entrant household per period is
\begin{equation}
 v_P=\sum_{j=0}^{J-1}(Q_P^\star)^j\gamma_P
     =(I-Q_P^\star)^{-1}\gamma_P.
 \label{eq:lifetime_measure}
\end{equation}
Let $\bar\varsigma=\prod_{\ell=1}^L\varsigma_\ell$ be survival from birth to
adult household formation. Define adult household-years per entrant and
lifetime next-generation households by
\begin{equation}
 \mathcal A(P)=\int_{\X}dv_P,
 \qquad
 R_0(P)=\nu\bar\varsigma\int_{\X}f_P(x)\,dv_P(x).
 \label{eq:renewal_objects}
\end{equation}
The first object determines the stationary entry rate needed per adult; the
second is the next-generation reproduction number. In an adult-normalized
stationary cross-section,
\begin{equation}
 E_0(P)=\frac{1}{\mathcal A(P)},
 \qquad
 B_0(P)=\frac{R_0(P)}{\mathcal A(P)},
 \qquad
 \boxed{\frac{B_0(P)}{E_0(P)}=R_0(P)}.
 \label{eq:BE_bridge}
\end{equation}
The three reproduction statistics used in this report are related but not
interchangeable. $\mathcal R_t$ is next-cohort reproduction in the
two-generation model, $R_0(P)$ is lifetime next-generation households per
lifecycle entrant, and the frozen demographic eigenvalue below is a one-period
asymptotic growth factor. Each equals one at its corresponding stationary
replacement point, but their values away from replacement need not coincide.

\begin{proposition}[Reproductive endpoint]
\label{prop:renewal_ss}
Suppose locally generated entrants enter with the same state distribution
$\gamma_P$ used in the renewal calculation. A nonzero closed-economy stationary
measure exists at $P^*$ if and only if
\begin{equation}
 R_0(P^*)=1.
 \label{eq:renewal_root}
\end{equation}
Conditional on this price, normalized composition is
$g^*=v_{P^*}/\mathcal A(P^*)$.
Housing clearing determines adult-household scale:
\begin{equation}
 S^*=\frac{K^s(P^*)}
 {\int_{\X}h(x;P^*)\,dg^*(x)}.
 \label{eq:renewal_scale}
\end{equation}
If $R_0$ is strictly decreasing, the reproductive price is unique.
\end{proposition}

The assumption on entrant states is substantive. Equation~\eqref{eq:BE_bridge}
is exact for the scalar renewal accounting used in the maintained model. It is
also an exact representation of an unrestricted demographic operator only if
descendants reproduce the wealth, earnings, and other state distribution used
for entrants. The present bequest motive does not itself deliver that result.

The stationary housing stock in equation~\eqref{eq:renewal_scale} is the
solution $K^s(P)$ to $\mathcal I(P,K^s(P);\zeta^S)=0$. Under a maintained fixed
stock it is the exogenous constant $\bar K$.

Under the same common entrant-state distribution assumption, an outside
entrant flow $M$ and local retention $\rho$ imply the stationary scale
\begin{equation}
 S^*=\frac{M}
 {E_0(P)-\rho B_0(P)},
 \qquad E_0(P)>\rho B_0(P).
 \label{eq:scalar_open_ss}
\end{equation}
Housing clearing then selects $P$. The closed economy is the transparent
theoretical benchmark. The open version is the empirically relevant
regularization when reproduction is below replacement and immigration is
positive. A national closure and a local housing market with migration are
therefore different equilibrium objects.

\section{Transition mechanics and interpretation}
\label{sec:transition_mechanics}

\subsection{Exact adult-household momentum}

Let $\sigma_t(x)=Q_t(x,\X)$ be the probability that an adult household remains
in the model next period, and let
\begin{equation}
 D_t=\int_{\X}[1-\sigma_t(x)]\,dG_t(x)
 \label{eq:adult_exits}
\end{equation}
be adult-household exit. Integrating equation
\eqref{eq:level_distribution_law} gives the exact accounting identity
\begin{equation}
 \boxed{S_{t+1}-S_t=e_{t+1}-D_t.}
 \label{eq:adult_momentum}
\end{equation}
Adult households grow when maturing locally born households plus immigrants
exceed exits. Current births do not enter equation~\eqref{eq:adult_momentum}
until those children mature. Thus period fertility, frozen-environment
reproduction, and current population growth are distinct objects.

This distinction resolves the apparent U.S. puzzle. Below-replacement
fertility today is consistent with a growing adult population when the current
age distribution reflects larger past birth cohorts or immigration. It is also
consistent with rising housing demand when those cohorts are moving into ages
with high ownership and space demand.

\subsection{Exact housing-demand decomposition}

Define expected next-period housing of a survivor by
\begin{equation}
 \begin{aligned}
 \int_{\X}h_{t+1}(x')Q_t(x,dx')
 &=\sigma_t(x)\bar h^S_{t+1}(x),\\
 e_{t+1}\bar h^{\rm new}_{t+1}
 &=\int_{\X}h_{t+1}(x)\,d\mathcal M^E_{t+1}(x)\\
 &\quad+\rho_tB_{t+1}\int_{\X}h_{t+1}(x)\,d\Gamma^B_{t+1}(x).
 \end{aligned}
 \label{eq:survivor_entrant_housing}
\end{equation}
Let $H_t^D=\int h_t(x)dG_t(x)$. The change in aggregate housing demand is
\begin{equation}
\begin{aligned}
 H_{t+1}^D-H_t^D
={}&\underbrace{\int\sigma_t(x)
 [\bar h^S_{t+1}(x)-h_t(x)]\,dG_t(x)}_{
 \text{aging, tenure transitions, and retention}}\\
&-\underbrace{\int[1-\sigma_t(x)]h_t(x)\,dG_t(x)}_{
 \text{housing released by exits}}
+\underbrace{e_{t+1}\bar h^{\rm new}_{t+1}}_{
 \text{new-household demand}}.
\end{aligned}
\label{eq:housing_decomposition}
\end{equation}
Equation~\eqref{eq:housing_decomposition} is an exact quantity identity. In
equilibrium its left side equals $K_{t+1}-K_t$, so its sign alone cannot predict
the next rent. A valid price test must also respect asset no-arbitrage. Hold the
continuation asset-price path $\mathcal P_{t+2:}$ fixed. For a candidate rent
$r$ at date $t+1$, equation~\eqref{eq:dynamic_user_cost} implies
\begin{equation}
 P_{t+1}(r)=
 \frac{r+P_{t+2}}{R_b+\delta_H+\tau^p_{t+1}}.
 \label{eq:candidate_asset_price}
\end{equation}
Let $\widehat h_{t+1}(x;r,P_{t+1}(r),\mathcal P_{t+2:})$ be the housing policy
under that candidate pair, and define
\begin{equation}
 \widehat H^D_{t+1}(r;\mathcal P_{t+2:})
 =\int_{\X}\widehat h_{t+1}
 \bigl(x;r,P_{t+1}(r),\mathcal P_{t+2:}\bigr)\,dG_{t+1}(x).
 \label{eq:old_rent_demand}
\end{equation}
Applying the same three terms in equation~\eqref{eq:housing_decomposition},
with every $h_{t+1}$ replaced by $\widehat h_{t+1}$, decomposes the shift in
this demand schedule into survivor aging and tenure, housing released by exits,
and new-household demand. The numerical packet computes the static-supply
analogue of this fixed-cost decomposition.

Because $K_{t+1}$ is predetermined by date-$t$ investment, if next-period
excess demand is strictly decreasing in its candidate rent, then
\begin{equation}
 r_{t+1}>r_t
 \quad\Longleftrightarrow\quad
 \widehat H^D_{t+1}(r_t;\mathcal P_{t+2:})>K_{t+1}.
 \label{eq:old_rent_price_test}
\end{equation}
This old-rent test is the full-model version of
Proposition~\ref{prop:toy_momentum}. It keeps the paper's
intergenerational mechanism visible: a large cohort can move from renting into
ownership, retain family-sized homes in old age, and shift demand outward even
as the entrant pipeline shrinks. The candidate asset price moves with the rent
through equation~\eqref{eq:candidate_asset_price}; the full asset-price change
still depends on the continuation path.

\subsection{Where the transition starts}

There need not be a single observed date at which a preference shock occurs.
The empirical transition should start from the first date at which the joint
state can be measured. Its initial condition must include the age--wealth--
tenure--parity distribution, the children already in the maturation pipeline,
and the housing stock. A smooth fertility wedge can then be inferred so that
the model matches age-specific birth hazards over time. This wedge is a compact
summary of omitted forces---childcare, work, partnership, norms, and other
costs---not an identified causal preference shock.

The appropriate comparison is therefore between two paths with the same
initial state. The baseline uses the inferred fertility-wedge sequence; the
no-decline counterfactual freezes that wedge at its initial value. Low
fertility can coexist with rising realized prices along the baseline because
of inherited momentum. But, holding the other paths fixed, it should lower
housing demand and prices relative to the no-decline counterfactual. This
relative statement is the causal discipline the levels alone cannot provide.

\subsection{Long-run policy incidence}

At a regular closed-economy root, any policy $g$ satisfies
\begin{equation}
 \frac{dP^*}{dg}=-\frac{R_{0g}(P^*;g)}{R_{0P}(P^*;g)}.
 \label{eq:root_comparative_static}
\end{equation}
If reproduction falls with the housing price, a policy that raises reproduction
at a fixed price supports a higher reproductive price. Relaxing young access
can have this effect because it changes tenure and fertility before children
are born. Its long-run population effect combines the induced price with
housing supply and per-household demand.

A pure old-release policy has different incidence. If it changes housing only
after fertility and child maturation, and does not enter earlier continuation
values, bequest utility, or taxes, then $R_{0g}=0$. It leaves the reproductive
price unchanged and accommodates more households by lowering housing per
adult. This neutrality is not automatic in the full model: anticipated
transaction costs and estate incentives can change earlier saving, tenure, and
fertility decisions. The quantitative model must compute, rather than assume,
the sign.

\section{A reproducible numerical transition}
\label{sec:numerical_transition}

The numerical exercise is deliberately small enough to audit. It demonstrates
the cohort accounting and policy timing before a full solution of the
lifecycle transition. It is not calibrated, not a forecast, and not a solved
asset-price equilibrium.

\subsection{Illustration}

One period is five years. Adult households occupy twelve age groups from
20--24 through 75--79. Let $x_{jt}$ be adult mass at age $j$, $q_{\ell t}$ a
four-cohort child-maturation queue measured in future adult-household
equivalents, $\sigma_j$ adult survival, $s_c$ child survival, and $\nu$ the
household-formation equivalents generated by one birth. The demographic law is
\begin{equation}
 x_{0,t+1}=s_cq_{3t},\qquad
 x_{j+1,t+1}=\sigma_jx_{jt},\qquad
 q_{0,t+1}=\nu\sum_j f_j(p_t,\theta_t,\omega_t)x_{jt},\qquad
 q_{\ell+1,t+1}=s_cq_{\ell t}.
 \label{eq:numerical_demography}
\end{equation}
Births occur at the end of period $t$ and enter $q_{0,t+1}$ at the next date.
The stages $q_0,\ldots,q_3$ represent ages 0--4 through 15--19; the last stage
enters $x_0$ at the following date. Thus adulthood occurs twenty years after
the end-of-period birth even though it is five date indices after $t$.
The fertility profile is concentrated around age 30. Its level falls smoothly
from a replacement environment as $\theta_t$ declines. The initial age
distribution has a prime-age bulge and smaller recent child cohorts. There is
no post-initial immigration.

At each date, a contemporaneous housing-cost index clears
\begin{equation}
 \bar H_s p_t^{\varepsilon}
 =\sum_j x_{jt}d_j(p_t;\theta_t,\omega_t,\tau_t).
 \label{eq:numerical_clearing}
\end{equation}
Housing demand rises through midlife and remains high at older ages. The
``frozen demographic growth factor'' is the dominant eigenvalue of the adult
aging, fertility, and maturation operator evaluated at current conditions.
It is below one when the current environment is below replacement.

\begin{figure}[ht!]
\centering
\includegraphics[width=0.94\textwidth]{../output/model/dynamic_population_housing_paper/transition_paths.pdf}
\caption{A generic fertility decline with inherited demographic momentum.
All paths start from the same adult and child distributions. The no-decline
path freezes the fertility schedule at its initial value. In the baseline,
housing cost and adult-household mass rise initially even after the frozen
demographic growth factor falls below one. Easier young access affects future
entry after children mature; lower old retention releases housing immediately.
All series are illustrative.}
\label{fig:numerical_paths}
\end{figure}

In the baseline, the frozen growth factor is below one after the initial date,
yet both adult households and the housing-cost index rise over several
five-year intervals. Adult households peak at $1.0325$ and the housing cost at
$1.0177$, both in year 25 relative to initial values of one. By year 50, the
indices have fallen to $0.9248$ and $0.9464$; by year 100 they are $0.7286$ and
$0.8159$. These magnitudes have no empirical content. The result is the sign
pattern: a declining endpoint can be preceded by decades of momentum.

The no-decline path starts from exactly the same inherited state. Its housing
cost is $1.0316$ in year 25 rather than $1.0177$ in the baseline. By year 50,
its adult-household and housing-cost indices are $0.9771$ and $0.9907$, versus
$0.9248$ and $0.9464$ in the baseline. Thus inherited momentum explains why
both realized paths can rise initially, while the gap between them isolates
the effect of the post-initial fertility-schedule decline in this illustration.

\begin{figure}[H]
\centering
\includegraphics[width=0.96\textwidth]{../output/model/dynamic_population_housing_paper/transition_accounting.pdf}
\caption{The inherited age profile and the exact fixed-price demand
decomposition. Early in the transition, maturing entrants and aging into
larger housing outweigh exits and the downward fertility schedule shift.
Later, exits dominate.}
\label{fig:numerical_accounting}
\end{figure}

The no-decline path keeps the value-of-children schedule at its initial level
while preserving the same inherited age distribution and child pipeline. The
easier-access experiment reduces the young-family access wedge after year
10. It raises fertility immediately but adult-household scale only after the
maturation delay. The lower-retention experiment reduces housing held at ages
60 and above. It lowers the clearing cost immediately and changes fertility
indirectly through that cost. Neither experiment is a welfare calculation.

The numerical packet verifies housing clearing, adult accounting, and the
demand decomposition to machine precision. Across all 80 nondegenerate
scenario transitions, the fixed-price demand sign predicts the next clearing
cost sign. The largest absolute housing residual, adult-accounting residual,
and decomposition residual are respectively $4.44\times10^{-16}$,
$2.64\times10^{-16}$, and $2.43\times10^{-16}$.

\subsection{What the illustration does not do}

The object $p_t$ in equation~\eqref{eq:numerical_clearing} is a current
housing-market clearing cost. It is not the forward-looking asset price $P_t$
in the lifecycle model. Solving the latter requires backward household
optimization and equation~\eqref{eq:dynamic_user_cost}. The exercise also
abstracts from income growth, interest-rate movements, construction, and
immigration after the initial date. It proves feasibility of the mechanism;
it does not explain the observed U.S. house-price path.

\section{Mapping to the maintained quantitative model}
\label{sec:quantitative_mapping}

\subsection{What is retained and what is new}

The intended dynamic model retains the maintained preferences and sequential
birth margins:
children shift equivalence-scaled needs toward housing; fertility is a
sequence of birth decisions; renters face a size cap; buyers face the
down-payment condition in equation~\eqref{eq:downpayment}; owners choose from a
housing ladder; and transaction costs and warm-glow bequests generate old-age
retention. It also makes one necessary architecture repair: independent child
stages replace the retained shared clock. Beyond that repair, the dynamic
closure adds the level distribution $G_t$, the aggregate child-maturation
pipeline needed to form future entrants, and time-indexed values and prices.
The retained benchmark parameters can therefore audit magnitudes and
identification, but they cannot be carried into this transition without
re-estimating the repaired fertility architecture.

This is not the one-shot completed-family-size model in the earlier circulated
draft. That version remains useful for local analytical results, but it cannot
be described as the same quantitative model as the retained equivalence-scale,
sequential-birth specification. The present document uses the latter as the
model that would be taken to the transition data.

The intergenerational emphasis is strengthened, not displaced. Young-access
policies have an immediate tenure and fertility effect and a delayed population
effect. Old-release policies have an immediate stock-allocation effect and an
indirect fertility effect through prices. Reporting both components reveals
which policies change family formation and which chiefly reallocate the
existing housing stock.

In a final paper, the two-generation model can serve as the theory section and
the lifecycle environment need be presented only once. The stationary benchmark
establishes the household and intergenerational allocation channels; one added
transition section then asks whether inherited cohorts reconcile low fertility
with rising housing costs and how much demographic feedback amplifies policy.

\subsection{Audited benchmark}

The best retained July 25 E5b parameter vector has current-contract loss
$385.8755$, asset price $0.815625$, and housing residual
$1.11\times10^{-6}$. It is an audited benchmark, not a promoted calibration:
the PSID rooms construction is under hold, and childlessness and old-wealth
dispersion remain large failures. Table~\ref{tab:selected_fit} reports a few
economically central rows; Appendix~\ref{app:fit} reports the complete target
system and every estimated or fixed parameter.

\begin{table}[ht!]
\centering
\caption{Selected audited fit of the retained benchmark}
\label{tab:selected_fit}
\small
\begin{tabular}{lrrl}
\toprule
Moment & Data & Model & Reading \\
\midrule
Completed fertility & 1.918 & 1.904 & Close \\
Family ownership gap & 0.168 & 0.166 & Close \\
Rooms gap, $3+$ versus 1--2 children & 0.368 & 0.372 & Close \\
Ownership, ages 30--55 & 0.575 & 0.564 & Close \\
Childless share & 0.188 & 0.080 & Large miss \\
Old wealth p90/p50 & 3.448 & 2.068 & Large miss \\
\bottomrule
\end{tabular}
\end{table}

The benchmark requires entrant flow $E_0=0.0617335$ per adult and produces
mature locally born flow $B_0=0.0511056$. Thus
\begin{equation}
 \frac{B_0}{E_0}=0.82784<1.
 \label{eq:audited_reproduction}
\end{equation}
Treating the $3+$ category with its reported top-code mean raises the diagnostic
ratio to $0.90319$, still below replacement. Across asset prices from 0.01 to
two times the benchmark, the maintained ratio ranges only from about 0.836 to
0.820; its local price elasticity is $-0.01067$. The maintained benchmark
therefore has no positive closed reproductive root over the audited range.

\begin{figure}[ht!]
\centering
\includegraphics[width=0.96\textwidth]{../output/model/closed_reproductive_closure_audit_20260811/maintained_closure_grid.png}
\caption{Reproductive and housing-scale maps at the retained benchmark. The
left panel shows that neither the model flow nor the $3+$ top-code diagnostic
reaches replacement over the audited price range. The right panel shows that
housing clearing nevertheless assigns stationary scale conditional on price.}
\label{fig:closure_audit}
\end{figure}

The conclusion is architecture-sensitive. With the retained parameters but an
independent maturation clock for each child, the audit gives $B_0/E_0=1.1179$;
after re-estimating that architecture in the August 5 diagnostic, it gives
$0.7852$. Neither exercise is a valid promoted counterfactual: the former does
not re-estimate behavior, and the latter uses a different fit. Their dispersion
shows why child timing must be fixed before estimating a dynamic closure.

\subsection{Identification consequence}

Stationary fertility levels and timing observed at one market price do not
identify the slope $R_{0P}$. The current target system contains strong moments
for levels, timing, tenure, and rooms, but no causal housing-cost-to-birth
hazard moment. The near-flat reproductive map in Figure~\ref{fig:closure_audit}
could therefore reflect either the model's true economic restriction or an
unidentified fertility-noise scale.

A quantitative transition needs a tenure-specific response moment that moves
housing cost or access while holding other child costs fixed. The estimand
should distinguish renters, prospective buyers, and incumbent owners because
the same asset-price movement raises rent and down payments but also raises
owner wealth. Without this slope evidence, the steady-state feedback and policy
amplification should be presented as theory, not as an estimated magnitude.

\section{Implementation and paper design}
\label{sec:implementation}

\subsection{A minimal full-transition algorithm}

The full transition can be implemented without redesigning the household
problem.
\begin{enumerate}
 \item Construct $G_0$ from the observed joint age, wealth, tenure, parity, and
 child-stage distribution. Construct $K_0$ from the housing stock. The initial
 child pipeline is required because it determines entrants over the first
 twenty years.
 \item Specify paths for the reduced-form fertility wedge, outside entry,
 construction shifters, taxes, and any policy. The baseline wedge matches
 age-specific birth hazards; the no-decline counterfactual freezes it.
 \item Guess paths for $P_t$ and $r_t$. Given a terminal value function, solve
 household problems backward using equation~\eqref{eq:bellman_path}.
 \item Push $G_t$ and the child pipeline forward one period using equations
 \eqref{eq:child_pipeline} and \eqref{eq:level_distribution_law}. Recompute births, entrants, exits,
 construction, and the period government budget.
 \item Update rents and asset prices to satisfy housing clearing and dynamic
 no-arbitrage. Iterate backward and forward until every date clears and the
 terminal distribution is close to its terminal stationary endpoint.
\end{enumerate}

The terminal endpoint can be open or closed. With the current benchmark it
should remain open in reported numerical exercises: no usable closed
reproductive root has been established over the audited price range. A changed
and re-estimated fertility architecture could alter that conclusion. A
period-by-period pension budget is also
required: the stationary pay-as-you-go formula cannot be reused when the age
distribution changes.

\subsection{Validation gates}

Before interpreting a computed transition, five gates should pass.
\begin{enumerate}
 \item Independent child clocks and the conversion of the $3+$ category must
 give consistent birth, maturation, and completed-fertility accounting.
 \item The initial age distribution and child pipeline must reproduce observed
 near-term adult-household growth before prices are used to fit anything.
 \item The baseline and no-decline paths must share $G_0$, $K_0$, migration,
 income, interest rates, and supply shifters; only the fertility wedge differs.
 \item Every period must report housing demand, supply, residuals, rents,
 asset prices, entrants, exits, the quantity identity in
 equation~\eqref{eq:housing_decomposition}, and the old-rent demand test in
 equation~\eqref{eq:old_rent_price_test}.
 \item Policy experiments must balance their fiscal accounts and separate the
 immediate household response from the delayed demographic-scale response.
\end{enumerate}

\section{Conclusion}

The model yields one clean separation. Reproductive balance selects a long-run
housing cost in the closed economy; housing clearing selects population scale.
The path between steady states is governed by inherited cohort composition,
tenure transitions, old-age retention, immigration, and durable supply. It can
therefore feature below-replacement reproduction, rising adult-household mass,
and rising housing costs at the same time.

The simple theory isolates the mechanism, and the numerical illustration
verifies its accounting. The lifecycle section provides a closed transition
architecture once its stated exogenous entry and policy paths are supplied.
What is not yet complete is a quantitative U.S. transition:
the retained benchmark has no closed reproductive root, child maturation is
architecture-sensitive, and the housing-cost elasticity of fertility is not
identified by the current moments. Those facts define a focused empirical and
computational agenda rather than a reason to abandon the closure.

The closest housing analogy is durable urban decline: as
\citet{glaeserGyourkoUrbanDecline2005} emphasize, demand and prices can move
well before the physical stock adjusts. The present mechanism adds the reverse
link from housing costs to family formation and the cohort structure through
which that link returns to future demand. The lifecycle fertility block is in
the tradition of \citet{sommerFertilityChoiceLife2016} and
\citet{doepkeBargainingBabiesTheory2019}; the housing side builds on tenure and
asset-market models such as \citet{hendersonModelHousingTenure1983} and
\citet{sommerEquilibriumEffectFundamentals2013}. The central synthesis is a
calendar-time renewal equilibrium that keeps these channels in one accounting
system; I do not claim that renewal theory itself is new.

\appendix

\section{Proofs}
\label{app:proofs}

\begin{proof}[Proof of Proposition~\ref{prop:toy_ss}]
The first-order conditions in equation~\eqref{eq:young_problem} give
$h=\alpha/p$ and
$n=\theta/[q+a(p+\omega)]$. Expenditures on adult housing and children are
$\alpha$ and $\theta$, so the stated income condition ensures positive
consumption. Reproductive stationarity requires $n=\bar n$, which gives the
price in equation~\eqref{eq:toy_ss}. It is positive exactly when
$\theta>\bar n(q+a\omega)$. In a steady state $Y=O$. Substituting young demand, old
demand, and $n=\bar n$ into market clearing and multiplying by $p$ gives the
stated population. It is increasing in $p$: its log derivative is
$(\varepsilon+1)/p-a\bar n/[\alpha+\varphi(\tau)+a\bar np]>0$. The comparative
statics follow directly.
\end{proof}

\begin{proof}[Proof of Proposition~\ref{prop:toy_momentum}]
At date $t$, total household population is $(1+\mu_t)Y_t$. At date $t+1$ it is
$(1+\mathcal R_t)Y_t$, so it rises exactly when $\mathcal R_t>\mu_t$. Evaluate next-period
housing demand at the current clearing price. Its change is
\[
 Y_t\{(\mathcal R_t-1)d_Y(p_t)+(1-\mu_t)d_O(p_t)\}.
\]
Supply is unchanged and increasing and demand is decreasing in price.
Therefore the next clearing price rises exactly when this expression is
positive, which is equation~\eqref{eq:price_momentum}. Combining the two strict
inequalities gives equation~\eqref{eq:joint_momentum}.
\end{proof}

\begin{proof}[Proof of Proposition~\ref{prop:renewal_ss}]
Let $e$ be stationary entrant flow. Iterating the survivor law gives
$G=e v_P$. The lifetime births generated by one entrant, multiplied by child
survival and the household-formation conversion, equal $R_0(P)$. In a
stationary child pipeline the mature flow therefore equals $eR_0(P)$: the
pipeline delays a constant flow but does not change it. A nonzero closed
stationary economy requires this flow to equal $e$, so $R_0(P)=1$.
Conversely, if $R_0(P)=1$, any positive multiple of $v_P$ reproduces its
entrant flow. Normalizing gives $g=v_P/\mathcal A(P)$, and housing
clearing fixes the multiple as in equation~\eqref{eq:renewal_scale}.
\end{proof}

\section{Complete benchmark target fit}
\label{app:fit}

\begin{table}[H]
\centering
\caption{Retained E5b benchmark: complete current-contract target system}
\label{tab:full_fit}
\scriptsize
\begin{tabularx}{\textwidth}{>{\raggedright\arraybackslash}Xrrrrr}
\toprule
Moment & Target & Model & Gap & Weight & Loss contribution \\
\midrule
Completed fertility & 1.918000 & 1.903573 & $-0.014427$ & 1425.739 & 0.297 \\
Childless share & 0.188000 & 0.080188 & $-0.107812$ & 17180.744 & 199.700 \\
Mean age at first birth & 25.310561 & 25.637433 & 0.326872 & 16.000 & 1.710 \\
Share of first births at age 30 or later & 0.270062 & 0.332334 & 0.062272 & 15625.000 & 60.590 \\
Rooms response at first birth (measurement hold) & 0.804944 & 0.657205 & $-0.147739$ & 35.735 & 0.780 \\
Rooms gap, $3+$ versus 1--2 children & 0.367700 & 0.372240 & 0.004540 & 2958.515 & 0.061 \\
Family ownership gap & 0.167662 & 0.165878 & $-0.001783$ & 14229.591 & 0.045 \\
Ownership rate, ages 30--55 & 0.575472 & 0.563758 & $-0.011714$ & 1207.846 & 0.166 \\
Mean occupied rooms, ages 18--85 & 5.779970 & 6.075505 & 0.295534 & 11.973 & 1.046 \\
Wealth / annual gross labor earnings & 6.873100 & 5.454754 & $-1.418346$ & 6.288 & 12.649 \\
Annual bequest flow / aggregate wealth & 0.008800 & 0.009078 & 0.000278 & 5165289.256 & 0.399 \\
Old wealth-to-income p90/p50, ages 76--84 & 3.448111 & 2.068370 & $-1.379741$ & 56.960 & 108.433 \\
\midrule
\multicolumn{5}{r}{Exact current-contract loss} & 385.875 \\
\bottomrule
\end{tabularx}
\end{table}

The table is an audit artifact, not a promoted calibration table. In
particular, the first-birth rooms row remains under a PSID measurement hold.
The unrounded row contributions sum exactly to the reported scalar; the small
difference visible after summing the printed entries is rounding.

\begin{table}[H]
\centering
\caption{Retained E5b benchmark: all free and externally fixed parameters}
\label{tab:full_parameters}
\scriptsize
\begin{tabularx}{\textwidth}{>{\raggedright\arraybackslash}X r c l}
\toprule
Parameter & Estimate & Search bound / restriction & Status \\
\midrule
Annual discount factor $\beta$ & 0.991097 & $[0.94,0.9995]$ & Estimated \\
Per-child expenditure shift $\Delta_1$ & 0.012163 & $[0,0.25]$ & Estimated \\
First-child expenditure jump $\Delta_0$ & 0.023896 & $[0,0.25]$ & Estimated \\
Value-of-child shifter $\psi$ & 0.585018 & $[-3,3]$ & Estimated \\
First-birth logit scale $\kappa_1$ & 3.552006 & $[0.02,50]$ & Estimated \\
Continuation logit scale $\kappa_C$ & 1.254006 & $[0.02,50]$ & Estimated \\
Owner-service premium $\chi$ & 1.038033 & $[0.1,5]$ & Estimated \\
Housing-supply scale $H_0$ & 8.162391 & $[0.2,80]$ & Estimated \\
Bequest intensity $\theta_0$ & 0.168435 & $[0,8]$ & Estimated \\
Bequest shift $\theta_1$ & 0.050046 & $[0.02,16]$ & Estimated; within 2\% of lower bound \\
Childless-renter expenditure share $\alpha_0$ & 0.733000 & CEX, 2019--2023 & Externally fixed \\
Tenure taste-shock scale & 0 & Author restriction & Externally fixed \\
Child dependence of bequest utility & 0 & Child-invariant bequest & Externally fixed \\
\bottomrule
\end{tabularx}
\end{table}

\section{Numerical illustration: parameters and replication}
\label{app:numerics}

\begin{table}[H]
\centering
\caption{Parameters of the generic transition illustration}
\label{tab:illustration_parameters}
\small
\begin{tabular}{lrlr}
\toprule
Object & Value & Object & Value \\
\midrule
Period length & 5 years & Adult ages & 20--79 \\
Maturation delay & 4 periods & Child survival per period & 0.992 \\
Household equivalents per birth $\nu$ & 0.50 & Initial housing-cost index & 1.00 \\
Housing-supply elasticity & 0.65 & Demand price elasticity & 1.00 \\
Child goods cost & 0.72 & Child space & 0.22 \\
Initial access wedge & 0.24 & Easier-access reduction & 0.12 \\
Initial value of children & 1.16 & Long-run value & 0.94 \\
Old-retention reduction & 18\% & Policy start & year 10 \\
\bottomrule
\end{tabular}
\end{table}

The pre-decline birth scale is chosen so the frozen demographic operator has
dominant eigenvalue one at a housing-cost index of one. The initial adult
distribution is a deterministic tilt of that stationary eigenvector, with a
prime-age bulge and modestly smaller recent child cohorts. The supply scale is
chosen so the initial market also clears at one. These normalizations create a
transparent sign experiment rather than an empirical calibration.

The two-panel steady-state figure is rebuilt from the repository root with
\begin{verbatim}
code/model/.venv/bin/python \
  code/model/tools/build_fertility_population_housing_transition_figures.py
\end{verbatim}
The complete numerical-transition packet is rebuilt with
\begin{verbatim}
code/model/.venv/bin/python \
  code/model/tools/build_dynamic_population_housing_paper_illustrations.py
\end{verbatim}
It writes aggregate paths, cohort distributions, the exact momentum-condition
grid, figures, parameters, and all accounting residuals to
\texttt{output/model/dynamic\_population\_housing\_paper/}.

\bibliographystyle{apalike}
\bibliography{main_bib,lit_review_extra}

\end{document}
